\documentclass[11pt]{article}
\usepackage{amsmath, amssymb, amsthm}
\usepackage[margin=1in]{geometry}
\usepackage{booktabs}
\usepackage{bm}

\newcommand{\pshed}{p^{\mathrm{shed}}}
\newcommand{\pserved}{p^{\mathrm{srv}}}
\newcommand{\pprev}{\pshed^{(\mathrm{prev})}}
\newcommand{\wprev}{w^{(\mathrm{prev})}}
\newcommand{\Top}{\mathcal{T}}
\newcommand{\Bot}{\mathcal{B}}

\title{Charnes--Cooper Reformulation of the Palma-Ratio Upper Level}
\author{FairLoadDelivery Package Documentation}
\date{May 2026}

\begin{document}
\maketitle

\section{Setup}

At bilevel iteration $k$, the lower-level multi-period MLD returns the load-shed
vector $\pprev \in \mathbb{R}^{m}$ and the Jacobian
$J = \partial \pshed / \partial w \in \mathbb{R}^{m\times m}$ via implicit
differentiation, where $m = T \cdot n$ stacks $n$ loads across $T$ periods. The
upper level chooses a weight update $\Delta w \in \mathbb{R}^{m}$ and uses the
first-order Taylor model
\begin{equation}\label{eq:taylor}
    \pshed(\Delta w) \;=\; \pprev + J\,\Delta w,
    \qquad
    \pserved(\Delta w) \;=\; \bm{p}^{d} - \pshed(\Delta w),
\end{equation}
subject to the trust region and weight box
\begin{align}
    |\Delta w_j| &\le \Delta, & j &= 1,\dots,m, \label{eq:trust}\\
    w_{\min} \le \wprev_j + \Delta w_j &\le w_{\max}, & j &\notin \mathrm{crit},\label{eq:wbox}
\end{align}
with $\Delta = 0.5$, $(w_{\min}, w_{\max}) = (1, 10)$, and critical loads
exempted from the upper bound.

\paragraph{Palma ratio (served).}
For a single period of $n$ loads with served vector $\pserved \in \mathbb{R}^{n}_{\ge 0}$,
let $\pserved_{(1)} \le \pserved_{(2)} \le \dots \le \pserved_{(n)}$ denote the
ascending order statistics, and define
\begin{equation}\label{eq:palma}
    \mathrm{Palma}(\pserved) \;=\;
    \frac{\sum_{i \in \Top} \pserved_{(i)}}{\sum_{i \in \Bot} \pserved_{(i)}}
    \;=\;
    \frac{\text{top 10\% served}}{\text{bottom 40\% served}},
\end{equation}
where $\Top = \{n - n_{\mathrm{top}}+1, \dots, n\}$ with $n_{\mathrm{top}} = \max(1,\lceil 0.1 n\rceil)$
and $\Bot = \{1,\dots, n_{\mathrm{bot}}\}$ with $n_{\mathrm{bot}} = \max(1,\lfloor 0.4 n\rfloor)$.

\section{Linearizing the Sort}\label{sec:sort}

Sorting is non-smooth, so we introduce a binary permutation matrix
$a^{(t)} \in \{0,1\}^{n\times n}$ per period $t$, with the doubly-stochastic
constraints
\begin{equation}\label{eq:dstoch}
    \sum_{j=1}^{n} a^{(t)}_{ij} = 1, \quad
    \sum_{i=1}^{n} a^{(t)}_{ij} = 1, \qquad i,j = 1,\dots,n.
\end{equation}
The product $u^{(t)}_{ij} = a^{(t)}_{ij}\,\pserved^{(t)}_{j}$ is bilinear in a
binary and a bounded-continuous variable. With
$0 \le \pserved^{(t)}_{j} \le p^{d}_{j}$ and $a^{(t)}_{ij} \in \{0,1\}$, the
exact McCormick envelope is
\begin{align}
    u^{(t)}_{ij} &\ge \pserved^{(t)}_{j} + a^{(t)}_{ij}\,p^{d}_{j} - p^{d}_{j}, \label{eq:mcc1}\\
    u^{(t)}_{ij} &\le a^{(t)}_{ij}\,p^{d}_{j}, \label{eq:mcc2}\\
    u^{(t)}_{ij} &\le \pserved^{(t)}_{j}, \label{eq:mcc3}\\
    u^{(t)}_{ij} &\ge 0.
\end{align}
For binary $\times$ non-negative bounded continuous, \eqref{eq:mcc1}--\eqref{eq:mcc3}
are tight at every binary realisation (they implement the disjunction exactly).
The sorted vector and the Palma sums become
\begin{equation}\label{eq:sorted}
    \mathrm{sorted}^{(t)}_{i} \;=\; \sum_{j=1}^{n} u^{(t)}_{ij},
    \qquad
    \mathrm{sorted}^{(t)}_{k} \le \mathrm{sorted}^{(t)}_{k+1}, \quad k = 1,\dots,n-1,
\end{equation}
\begin{equation}\label{eq:sums}
    \mathrm{top\_sum}^{(t)} \;=\; \sum_{i\in\Top} \mathrm{sorted}^{(t)}_{i},
    \qquad
    \mathrm{bot\_sum}^{(t)} \;=\; \sum_{i\in\Bot} \mathrm{sorted}^{(t)}_{i}.
\end{equation}

\section{Weak Charnes--Cooper (MIQCP)}\label{sec:weak}

The Palma ratio in \eqref{eq:palma} is a linear-fractional functional of the
sorted served vector. The classical Charnes--Cooper transformation
\cite{CharnesCooper1962} introduces a scaling $\sigma > 0$ that normalises the
denominator:
\begin{equation}\label{eq:cc-norm}
    \sigma^{(t)} \cdot \mathrm{bot\_sum}^{(t)} \;=\; 1,
    \qquad \sigma^{(t)} \ge \sigma_{\min} > 0.
\end{equation}
Under \eqref{eq:cc-norm}, minimising the Palma ratio is equivalent to
\begin{equation}\label{eq:weak-obj}
    \min \quad \sum_{t=1}^{T} \lambda_{t}\,\sigma^{(t)}\,\mathrm{top\_sum}^{(t)}.
\end{equation}
We refer to this as the \emph{weak} Charnes--Cooper formulation because
$\sigma$ is multiplied in at only two places (constraint \eqref{eq:cc-norm} and
objective \eqref{eq:weak-obj}). The remaining decision variables ($\Delta w$,
$a^{(t)}$, $u^{(t)}$) keep their original meaning. The full upper-level problem
is
\begin{align}
    \min_{\Delta w,\,a,\,u,\,\sigma} \quad &
        \sum_{t=1}^{T} \lambda_{t}\,\sigma^{(t)}\,\mathrm{top\_sum}^{(t)}
        \label{eq:weak-full}\\
    \text{s.t.} \quad &
        \eqref{eq:taylor},\ \eqref{eq:trust},\ \eqref{eq:wbox},\nonumber\\
        & \varepsilon \le \pshed_{j} \le p^{d}_{j}, \quad j = 1,\dots,m,\nonumber\\
        & \eqref{eq:dstoch},\ \eqref{eq:mcc1}\text{--}\eqref{eq:mcc3},\ \eqref{eq:sorted},\ \eqref{eq:sums},\nonumber\\
        & \sigma^{(t)}\,\mathrm{bot\_sum}^{(t)} = 1, \quad t = 1,\dots,T.\nonumber
\end{align}
The bilinear products $\sigma^{(t)}\,\mathrm{bot\_sum}^{(t)}$ and
$\sigma^{(t)}\,\mathrm{top\_sum}^{(t)}$ make this a non-convex MIQCP. Gurobi
solves it with \texttt{NonConvex=2}, but its LP relaxation is loose (Gurobi
frequently gets stuck at the do-nothing incumbent) and runtime degrades sharply
in $T\cdot n$.

\section{Formal Charnes--Cooper (MILP)}\label{sec:formal}

The standard Charnes--Cooper trick recovers an LP only when \emph{every}
original variable that appears multiplicatively with the denominator is
rescaled by $\sigma$. We apply this rescaling load-by-load, period-by-period.
Let $t(j) = \lfloor (j-1)/n \rfloor + 1$ denote the period of global index $j$.
For each period $t$ we introduce
\begin{equation}\label{eq:rescale}
    \Delta w^{z}_{j} \;=\; \sigma^{(t(j))}\,\Delta w_{j},
    \qquad
    \pserved^{z}_{j} \;=\; \sigma^{(t(j))}\,\pserved_{j},
    \qquad
    u^{z,(t)}_{ij} \;=\; \sigma^{(t)}\,u^{(t)}_{ij}.
\end{equation}

\subsection{Block-diagonality of the Jacobian}

Because the multi-period lower-level MLD is fully separable across periods
(weights, dispatch, and shed in period $t$ depend only on period-$t$ data),
$J$ is block-diagonal with $n\times n$ blocks: $J_{jk} = 0$ whenever
$t(j) \ne t(k)$. This is essential — it lets us rescale period $t(j)$'s
weight perturbations by $\sigma^{(t(j))}$ without coupling to other periods'
$\sigma$. A runtime guard (in \texttt{palma\_ratio\_minimization\_formal\_cc})
checks $\max_{t(j)\ne t(k)} |J_{jk}|$ against a tolerance and warns if the
assumption is violated.

\subsection{Rescaled Taylor and sums}

Multiplying \eqref{eq:taylor} by $\sigma^{(t(j))}$ and using block-diagonality,
\begin{equation}\label{eq:taylor-z}
    \pserved^{z}_{j}
    \;=\; \sigma^{(t(j))}\bigl(p^{d}_{j} - \pprev_{j}\bigr)
        - \sum_{k=1}^{m} J_{jk}\,\Delta w^{z}_{k}.
\end{equation}
The doubly-stochastic constraint \eqref{eq:dstoch} is unchanged (it's already
linear in the binaries). The normalising equation \eqref{eq:cc-norm} becomes
\begin{equation}\label{eq:cc-norm-z}
    \sum_{i\in\Bot}\,\sum_{j=1}^{n} u^{z,(t)}_{ij} \;=\; 1, \qquad t = 1,\dots,T,
\end{equation}
and the objective \eqref{eq:weak-obj} becomes
\begin{equation}\label{eq:formal-obj}
    \min \quad
    \sum_{t=1}^{T} \lambda_{t} \sum_{i\in\Top}\,\sum_{j=1}^{n} u^{z,(t)}_{ij}.
\end{equation}
Both are linear in the rescaled variables.

\subsection{Rescaled box constraints}

Multiplying every RHS constant by $\sigma^{(t)}$ preserves the constraint under
the substitution \eqref{eq:rescale}. The trust region \eqref{eq:trust} becomes
\begin{equation}\label{eq:trust-z}
    -\Delta\,\sigma^{(t(j))} \;\le\; \Delta w^{z}_{j} \;\le\; \Delta\,\sigma^{(t(j))},
\end{equation}
the weight box \eqref{eq:wbox} becomes
\begin{equation}\label{eq:wbox-z}
    w_{\min}\,\sigma^{(t(j))} \;\le\; \wprev_{j}\,\sigma^{(t(j))} + \Delta w^{z}_{j} \;\le\; w_{\max}\,\sigma^{(t(j))},
\end{equation}
and the pshed bounds $\varepsilon \le \pshed_{j} \le p^{d}_{j}$ — equivalently
$0 \le \pserved_{j} \le p^{d}_{j} - \varepsilon$ — become
\begin{equation}\label{eq:pbox-z}
    0 \;\le\; \pserved^{z}_{j} \;\le\; (p^{d}_{j} - \varepsilon)\,\sigma^{(t(j))}.
\end{equation}
All of \eqref{eq:trust-z}--\eqref{eq:pbox-z} are linear in
$(\Delta w^{z}, \sigma)$.

\subsection{The bilinear $u^{z} = a \cdot \pserved^{z}$ via indicator constraints}

The only remaining bilinearity is the product $u^{z,(t)}_{ij} = a^{(t)}_{ij}\,\pserved^{z}_{t(\cdot),j}$.
Because $a^{(t)}_{ij}$ is binary, the product can be encoded \emph{exactly} as
a disjunction:
\begin{align}
    a^{(t)}_{ij} = 1 \;&\Longrightarrow\; u^{z,(t)}_{ij} = \pserved^{z}_{(t-1)n+j}, \label{eq:ind1}\\
    a^{(t)}_{ij} = 0 \;&\Longrightarrow\; u^{z,(t)}_{ij} = 0, \label{eq:ind2}\\
    u^{z,(t)}_{ij} &\ge 0.\nonumber
\end{align}
JuMP/Gurobi natively support indicator constraints of the form
\texttt{a => \{lin\_constraint\}}; Gurobi enforces them through its branching
tree (SOS1-style logic) without injecting any Big-M into the model. This is
the second key step that makes the formal CC variant a clean MILP. The earlier
McCormick implementation required a finite upper bound $\sigma_{\max}^{(t)}$ to
bound $\pserved^{z}_{j} \le (p^{d}_{j} - \varepsilon)\,\sigma_{\max}^{(t)}$;
choosing $\sigma_{\max}^{(t)} = 10 / \mathrm{bot\_sum}_{\mathrm{prev}}^{(t)}$
created an artificial \texttt{INFEASIBLE} mode when bilevel iterations pushed
$\mathrm{bot\_sum}_{\mathrm{prev}}^{(t)} \to 0$. Indicator constraints remove
the need for any upper bound on $\sigma^{(t)}$, matching the formal CC math
($\sigma^{(t)} = 1/\mathrm{bot\_sum}^{(t)}$ may grow unboundedly).

\subsection{The full MILP}

\begin{align}
    \min_{\substack{\Delta w^{z},\,a,\\u^{z},\,\sigma}}
    \quad &
    \sum_{t=1}^{T} \lambda_{t} \sum_{i\in\Top}\,\sum_{j=1}^{n} u^{z,(t)}_{ij}
    \label{eq:formal-full}\\
    \text{s.t.}\quad
    & \pserved^{z}_{j}
        = \sigma^{(t(j))}(p^{d}_{j} - \pprev_{j}) - \sum_{k=1}^{m} J_{jk}\,\Delta w^{z}_{k},
        & j &= 1,\dots,m,\nonumber\\
    & 0 \le \pserved^{z}_{j} \le (p^{d}_{j} - \varepsilon)\,\sigma^{(t(j))},
        & j &= 1,\dots,m,\nonumber\\
    & \sum_{j=1}^{n} a^{(t)}_{ij} = 1,\quad \sum_{i=1}^{n} a^{(t)}_{ij} = 1,
        & t,i,j,\nonumber\\
    & a^{(t)}_{ij} = 1 \Rightarrow u^{z,(t)}_{ij} = \pserved^{z}_{(t-1)n + j},
        & t,i,j,\nonumber\\
    & a^{(t)}_{ij} = 0 \Rightarrow u^{z,(t)}_{ij} = 0,
        & t,i,j,\nonumber\\
    & \sum_{j=1}^{n} u^{z,(t)}_{kj} \le \sum_{j=1}^{n} u^{z,(t)}_{k+1,j},
        & t,\, k &= 1,\dots,n-1,\nonumber\\
    & \sum_{i\in\Bot}\,\sum_{j=1}^{n} u^{z,(t)}_{ij} = 1,
        & t &= 1,\dots,T,\nonumber\\
    & -\Delta\,\sigma^{(t(j))} \le \Delta w^{z}_{j} \le \Delta\,\sigma^{(t(j))},
        & j &= 1,\dots,m,\nonumber\\
    & w_{\min}\,\sigma^{(t(j))}
        \le \wprev_{j}\,\sigma^{(t(j))} + \Delta w^{z}_{j}
        \le w_{\max}\,\sigma^{(t(j))},
        & j &\notin \mathrm{crit},\nonumber\\
    & \sigma^{(t)} \ge \sigma_{\min} > 0,
        & t &= 1,\dots,T,\nonumber\\
    & a^{(t)}_{ij} \in \{0,1\},\quad u^{z,(t)}_{ij} \ge 0.\nonumber
\end{align}
Every relation is linear in $(\Delta w^{z}, \pserved^{z}, u^{z}, \sigma)$ aside
from the binaries $a^{(t)}_{ij}$ and the indicators in \eqref{eq:ind1}--\eqref{eq:ind2},
so the problem is a (pure) MILP under Gurobi's native indicator support.
No \texttt{NonConvex=2} is needed.

\subsection{Recovery of original-space variables}

Given an optimal $(\Delta w^{z*}, \sigma^{*})$, the original-space weight
update is
\begin{equation}\label{eq:recover}
    \Delta w_{j}^{*} \;=\; \frac{\Delta w^{z*}_{j}}{\sigma^{*(t(j))}},
    \qquad
    \pshed^{*} \;=\; \pprev + J\,\Delta w^{*},
    \qquad
    w^{*} \;=\; \wprev + \Delta w^{*}.
\end{equation}
$\sigma^{*(t)} \ge \sigma_{\min} > 0$ guarantees the division in
\eqref{eq:recover} is well-defined. Because the rescaling \eqref{eq:rescale}
preserves the feasible set of the original Palma problem (the rescaled
constraints are scalar multiples by the positive $\sigma^{(t)}$ of the
originals), $(w^{*}, \pshed^{*})$ is feasible and optimal for the
weak-CC problem \eqref{eq:weak-full}.

\section{Summary of the Two Variants}

\begin{table}[h]
\centering
\small
\begin{tabular}{@{}lll@{}}
\toprule
 & \textbf{Weak CC \eqref{eq:weak-full}} & \textbf{Formal CC \eqref{eq:formal-full}} \\
\midrule
Class & MIQCP (non-convex) & MILP \\
Bilinearities kept & $\sigma\cdot\mathrm{bot\_sum},\ \sigma\cdot\mathrm{top\_sum},\ a\cdot \pserved$ & none (indicators replace $a\cdot \pserved^{z}$) \\
Gurobi setting & \texttt{NonConvex=2} & default (MILP) \\
$\sigma$ upper bound & not needed (but LP-relax loose) & not needed (indicators, no Big-M) \\
Rescaled variables & --- & $\Delta w,\ \pserved,\ u$ (i.e.\ multiplied by $\sigma$) \\
Recovers original $\Delta w$ & directly & via $\Delta w = \Delta w^{z} / \sigma$ \\
Block-diagonal $J$ required & no & yes (warned on violation) \\
Speed (case6, $T{=}8$) & baseline & $\sim 2.3\times$ faster \\
\bottomrule
\end{tabular}
\caption{Weak vs.\ formal Charnes--Cooper variants of the Palma upper level.
The implementation defaults to formal CC; the weak-CC path is retained as a
fallback via the \texttt{use\_weak\_cc=true} flag.}
\end{table}

\begin{thebibliography}{1}
\bibitem{CharnesCooper1962}
A.~Charnes and W.~W. Cooper.
\newblock Programming with linear fractional functionals.
\newblock {\em Naval Research Logistics Quarterly}, 9(3--4):181--186, 1962.
\end{thebibliography}

\end{document}
